\documentclass[11pt]{article}
\usepackage{setspace}
\usepackage[top=1in, bottom=1in, left=1in, right=1in]{geometry}
\usepackage{hyperref}
\usepackage{enumitem}
\usepackage[normalem]{ulem}
\usepackage{changepage}
\usepackage{xcolor}
\usepackage{ulem}

% =========================
% Custom macros for coauthor comments
% =========================
\newcommand{\sst}[1]{\textcolor{green}{#1}}
\newcommand{\tk}[1]{\textcolor{orange}{#1}}
\newcommand{\mw}[1]{\textcolor{magenta}{#1}}
% ===========================================================================

\onehalfspacing

% Title information
\title{Preferences on Tax Enforcement and Tax Administration:\\Draft of Survey Questions}
\author{Tim Kircali, Stefanie Stantcheva, Matthias Weber}
\date{\today} % Use \date{} for no date

% Bibliography package
\usepackage[backend=biber, style=apa]{biblatex}
\addbibresource{references.bib} % Link to the bibliography file

\begin{document}
\maketitle

\noindent We indicate with ”/” and ”-” whether questions have discrete choice options (e.g. yes / no) or continuous choice options (e.g. strongly oppose - strongly support). 
\section*{Page for the Confirmation of Participant Consent}
\section*{Demographic Background}
\begin{enumerate}
    \item What is your gender? \\ \textit{(Male/Female/Other)} 
    \item How old are you? \\ \textit{(-)[nr. between 18 - 100]}
    \item What is your zipcode? \\ \textit{(-)[a valid zip code]}
    \item What is your highest level of education? \\ \textit{(No High School degree / High School degree/ College Degree / Master’s or Professional (JD, MD, MBA) Degree / Doctoral Degree)}
    \item What is your ethnicity? \\ \textit{(European American-White / African American-Black / Hispanic-Latino / Asian-Asian American / Mixed race or other)}
    \item Do you live with your partner (if you have one)? \\ \textit{(Yes/No or I do not have a partner)}
    \item How many people are you in your household? (a household only includes members of family or your dependants but not roommates) \\ \textit{(1/2/3/4/5 or more)}
    \item How many children below the age of 14 live with you? \\ \textit{(-)}
    \item Are you a homeowner or a tenant? \\ \textit{(Homeowner/Tenant/Both/Other)}
    \item What is your employment status? \\ \textit{(Full-time employed/ Part-time employed/ Self-employed/ Student/ Retired/ Unemployed (searching for a job)/ Inactive (no job and not searching for one))}
    \item What was your annual household income in 2024 [before withholding tax]?\\ 
    \textit{(Less than [Q1] / Between [Q1] and [Q2] / Between [Q2] and [Q3] / Between [Q3] and [Q4] / More than [Q4] / Prefer not to say)}
    \item What are your current sources of income? (Select all that apply)\\
    \textit{(Salary or wages from employment / Self-employment or freelance income / Rental income / Investment income (e.g., dividends, capital gains) / Cash-based or informal income / Other)}
    \item What is the estimated total value of your assets — or your household’s assets if you are married? Please include everything you own (e.g., home, car, savings, investments), minus any outstanding debts (such as mortgages, car loans, or credit card balances). For example: if your home is worth \$300,000 and you still owe \$200,000 on your mortgage, the net value of your home is \$100,000.\\
    \textit{(Less than [Q1] / Between [Q1] and [Q2] / Between [Q2] and [Q3] / Between [Q3] and [Q4] / More than [Q4] / Prefer not to say)}
\end{enumerate}
\section*{Political Leaning}
\tk{Q to SS: Which of the following questions do we need? \textbf{We suggest to remove question 14.}}
\begin{enumerate}[resume]
    \stepcounter{enumi}
    \item[\sout{\theenumi.}] \textcolor{gray}{To what extend are you interested in politics?\\ \textit{(1: Not interested at all, 2: Slightly interested, 3: Moderately interested, 4: Interested, 5: Very interested)}}
    \item Did you vote in the last presidential elections? \\\textit{(Yes/No)} 
    \item Who did you support in the last presidential elections? \\ \textit{(Donald Trump/ Kamela Harris/ Other)}
    \item On economic policy matters, where do you see yourself on the liberal/conservative spectrum? \\\textit{(1: Very liberal, 2: Liberal, 3: Neither liberal nor conservative, 4: Conservative, 5: Very Conservative)}
    \item What do you consider your political affiliation, as of today? \\
    \textit{(Republican / Democrat / Independent / Other / Non-Affiliated)} \\
    \item (If “Other” to 16) Please specify your political affiliation: \textit{Text box}
\end{enumerate}
\section*{Government Views \& Trust}
\tk{Q to SS: Which of the following questions do we need? Can we remove 1-2 questions? \textbf{We suggest to remove question 20 and 22.}}\\
\underline{Trust}
\begin{enumerate}[resume]
    \stepcounter{enumi}
    \item[\sout{\theenumi.}] \textcolor{gray}{Do you agree or disagree with the following statement: “Most people can be trusted.” \\
    \textit{(1: Strongly disagree, 2: Somewhat disagree, 3: Neither disagree nor agree, 4: Somewhat Agree, 5: Strongly agree)}}
\end{enumerate}
\begin{adjustwidth}{0.75em}{0em}The following questions ask about your general views on the federal government. When answering, \uline{please think broadly about the past 10 years and your expectations for the next 10 years, and try to set aside the actions of any specific administration.}\end{adjustwidth}

\begin{enumerate}[resume]
    \item On average how often do you think you can trust the federal government to do what is right? \\ \textit{(1: Almost never, 2: Not very often, 3: Sometimes, 4: A lot of the time - 5: Almost always)} 
\end{enumerate}
\underline{Desired Involvement of Government}
\begin{enumerate}[resume]
    \stepcounter{enumi}
    \item[\sout{\theenumi.}] \textcolor{gray}{Where would you rate yourself on a scale from 1 to 10, where 1 means you think the government should do only those things necessary to provide the most basic government functions, and 10 means you think the government should take active steps in many areas trying to improve the lives of its citizens? You may use any number from 1 to 10. \\\textit{(1: government should be very inactive - 10: government should be very active)}}
\end{enumerate}
\begin{enumerate}[resume]
    \item Do you think the government should in general be more or less involved in providing public services and functions than it is now?
    \textit{(1: Much less involved, 2: Less involved, 3: About the same, 4: More involved, 5: Much more involved)}

\end{enumerate}
\underline{Quality of Government Services}
\begin{enumerate}[resume]
    \item Now think about the quality of the services that the government provides to citizens (such as registration, licensing, healthcare, or social support)? Consider how easy these services are to use, how quickly they are delivered, and how affordable they are. How would you rate government services in your country? \\
    \textit{(1: Very poor, 2: Poor, 3: Average, 4: Good, 5: Very good)}
\end{enumerate}
\section*{First Order Concerns}
\tk{Q to SS: Open-ended questions help us to better understand people's concerns by avoiding any restrictions or priming of respondents regarding certain considerations. The following questions could be interesting. However, we are hesitant to include open-ended questions because we are unsure whether the results would be worth the effort of analysing them. Do you think it is important to include open-ended questions such as the following? \textbf{Our suggestion is to not include question 25 and 26.}}
\begin{enumerate}[resume]
    \stepcounter{enumi}
    \item[\sout{\theenumi.}] \textcolor{gray}{What are your main considerations, whether positive or negative, regarding the process by which U.S. taxpayers file their taxes? Are there any problems that the government should address?
    \textit{Open Ended}}
    \stepcounter{enumi}
    \item[\sout{\theenumi.}] \textcolor{gray}{Now if you think about tax evasion in the U.S, what are your main considerations that come to your mind? Are there any problems the government should address?
    \textit{Open Ended}}
\end{enumerate}
\section*{Personal Exposure to Taxes}
\tk{Q to SS: Potentially, we could remove question 28 or/and 29. Within consideration questions we already ask whether tax filing costs are high and later we ask whether they respondents support prepopulated tax returns. \cite{benzarti2024rising} also ask for WTP for prepopulated tax returns, therefore question 30 could be an interesting comparison. \textbf{Still, I suggest to remove question 29 and 30 to save space.}} \\
\underline{Exposure to Tax Filing Costs}
\begin{enumerate}[resume]
    \item People file their tax returns in different ways. Some prepare and submit them on their own, while others rely on help from friends, relatives, or professionals. How do you usually file your taxes?\\ 
    \textit{(I prepare and file them myself / I prepare them together with help from a friend or relative /  A friend or relative prepares and files them for me / I prepare them together with a professional / A professional prepares and files them for me / I am not filing my taxes, and neither is anyone else / Other)}
    \item Some peoples' tax returns are filed using tax software, while those of others are filed by hand. How are your tax returns filed, by hand or by using tax software? \\ 
    \textit{(By hand / By using tax software / Both, by hand and by using tax software / I don't know / I am not filing my taxes, and neither is anyone else)}
    \stepcounter{enumi}
    \item[\sout{\theenumi.}] \textcolor{gray}{How tedious do you think that the task of filing taxes is? \\ \textit{(1: Not tedious at all, 2: Slightly tedious, 3: Moderately tedious, 4: Tedious, 5: Very tedious)}} 
    \stepcounter{enumi}
    \item[\sout{\theenumi.}] \textcolor{gray}{What is the maximum amount you would be willing to pay for a tax-filing service that automatically prepares your tax return, requiring only your confirmation before submission and no effort on your part to collect the necessary documents? \\
    \textit{(Slider) in dollars}}
\end{enumerate}
\underline{Exposure to Tax Enforcement}
\vspace{0.5em} 
\begin{adjustwidth}{0.75em}{0em}For the next three questions, please evaluate how likely you think it is that other people will do certain things when filing their taxes.\end{adjustwidth}
\begin{enumerate}[resume]
    \item How likely do you think it is that people you know personally unintentionally make mistakes in their tax returns ? \\
    \textit{(1: Very unlikely, 2: Unlikely, 3: Neither likely nor unlikely, 4: Likely, 5: Very likely
 / Prefer not to say)}
    \item Some people intentionally report minor details incorrectly to reduce their tax burden. Thinking about people you know personally, how likely do you think it is that they intentionally report minor details incorrectly?\\
    \textit{(1: Very unlikely, 2: Unlikely, 3: Neither likely nor unlikely, 4: Likely, 5: Very likely
 / Prefer not to say)}
    \item Others may go further and deliberately conceal significant parts of their income to avoid paying higher taxes. Among people you know personally, how likely do you think it is that they hide significant parts of their income to avoid paying higher taxes?\\
    \textit{(1: Very unlikely, 2: Unlikely, 3: Neither likely nor unlikely, 4: Likely, 5: Very likely
 / Prefer not to say)}
\end{enumerate}
\section*{Knowledge}
\tk{Q to SS: We are not sure whether we should include a special introduction to the knowledge questions - or potentially even incentivize the questions. In your recent health care paper you used the following paragraph. Should we include it too? \textbf{We suggest not including it to save space.}}
\\
\begin{adjustwidth}{0.75em}{0em}
\textcolor{gray}{As you probably know, the government and researchers gather a lot of statistical information about the economy. We are interested in learning whether this information finds its way to the general public. The next set of questions is about some economic policies in the United States. These are questions for which there are right or wrong answers.}

\textcolor{gray}{In order for your answers to be most helpful to us, it is really important that you answer these questions as accurately as you can and without consulting any external sources. Although you may find some questions difficult, it is very important for our research that you try your best. Thank you very much!}
\end{adjustwidth}
\vspace{0.5em}
\begin{adjustwidth}{0.75em}{0em}Please carefully read the following explanations and keep it in mind when answering the next question.\end{adjustwidth} 
\vspace{0.5em}
\begin{adjustwidth}{1.25em}{1.25em} 
\textbf{Explanation Tax Evasion:} Some individuals and companies do not pay their taxes fully. They may hide income or assets from the tax authorities, or they may misreport information on their tax returns. This is called tax evasion. 
\end{adjustwidth}
\begin{adjustwidth}{1.25em}{1.25em} \textbf{Explanation Tax Filing Costs:} Tax filing costs are the costs incurred when filing taxes. These include expenses for software license fees, expenses for paying someone else to file tax returns, but also the time cost when filing the taxes. 
\end{adjustwidth}
\vspace{0.5em}
\begin{enumerate}[resume]
    \item Please think about tax evasion. If you had to take a guess, out of each 100 dollars of total tax revenue how many dollars do you think is lost due to tax evasion? \\
    \textit{(Slider 0-100) dollars} 
\end{enumerate}
\begin{enumerate}[resume]
    \item Please think about tax filing costs. What do you think, on average, how much time and how many dollars do US individual taxpayers spend to file their taxes? \\
    \textit{Time: (Slider) hours} \\
    \textit{Expenses: (Slider) dollars}
\end{enumerate}
\vspace{0.5em}
\tk{Q to SS: \textbf{We suggest to not use question 36 and 37.} The insights would be interesting but: 1. The questions are complex. 2. We don't provide respondents with the relevant information. 3. For us, in fact, it would be difficult to validate whether their responses are correct. 4. We could save response time by not including all these subquestions. 5. We might have some problems with priming. \textbf{What do you think about it? Are there any other knowledge questions that might be very important to include?}}

\begin{enumerate}[resume]
    \stepcounter{enumi}
    \item[\sout{\theenumi.}] \textcolor{gray}{Some groups of taxpayers may contribute differently to the missing tax revenue. If you had to take a guess, out of each 100 dollars of total tax revenue, how many dollars do you think are lost due to tax evasion within each of the following groups? (In red you can see the level of tax evasion you indicated before for the overall population.)}
    \textcolor{gray}{
    \begin{enumerate}[label=\alph*)]
        \item Large businesses \textit{(Slider 0-100) dollars} 
        \item Small businesses \textit{(Slider 0-100) dollars} 
        \item High-income households \textit{(Slider 0-100) dollars} 
        \item Middle-income households \textit{(Slider 0-100) dollars} 
        \item Low-income households \textit{(Slider 0-100) dollars} 
    \end{enumerate}}
\end{enumerate}
\begin{enumerate}[resume]
    \stepcounter{enumi}
    \item[\sout{\theenumi.}] \textcolor{gray}{The tax authority checks if tax returns are correct. For that they use various data sources. While they have direct access to some types of information, other data is only accessible under specific conditions, such as a strong suspicion of tax fraud.\\
    What do you think, to how much of the following types of tax-relevant information does your country’s tax authority (the IRS) have direct and automatic access — without needing a special request or court order? \\ 
    Please answer using a scale from 1 to 10, where 1 is no access at all and 10 is full access.}
    \textcolor{gray}{
    \begin{enumerate}[label=\alph*)]
        \item Employment income (e.g., wages, salaries, and employer-reported income via W-2 or 1099 forms) \\ \textit{(1: No access, 2: Between no and limited access, 3: Limited access, 4: Between limited and unlimited  access, 5: Unlimited access)} 
        \item Self-employment income (e.g., income from freelance work or own business not reported by a third-party) \\ \textit{(1: No access, 2: Between no and limited access, 3: Limited access, 4: Between limited and unlimited  access, 5: Unlimited access)} 
        \item Investment and capital income (e.g., dividends, interest payments, or rental income reported by banks or platforms) \\ \textit{(1: No access, 2: Between no and limited access, 3: Limited access, 4: Between limited and unlimited  access, 5: Unlimited access)}
        \item Health insurance data (e.g., information about coverage status or premium amounts) \\ \textit{(1: No access, 2: Between no and limited access, 3: Limited access, 4: Between limited and unlimited  access, 5: Unlimited access)}
        \item Bank account data (e.g., account balances or individual transactions) \\ \textit{(1: No access, 2: Between no and limited access, 3: Limited access, 4: Between limited and unlimited  access, 5: Unlimited access)}
        \item Payment data (e.g., credit card transactions, mobile payments, or digital wallet activity) \\ \textit{(1: No access, 2: Between no and limited access, 3: Limited access, 4: Between limited and unlimited  access, 5: Unlimited access)}
    \end{enumerate}}
\end{enumerate}
\section*{Experimental Treatments}
\textit{This section only applies to respondents outside the control treatment. We introduce and then display the experimental video treatments. Afterwards, we ask a few questions to validate that people were actually watching the video. These questions should be kept very simple, only checking whether they watched the video. That there will be such questions is already mentioned in the introductions to the videos.}
\begin{enumerate}[resume]
    \item Introduction to video and validation questions
    \begin{itemize}
        \item Evasion: Revenue
        \begin{itemize}
            \item Introduction: tbd
            \item Validation Questions: tbd
        \end{itemize}
        \item Evasion: Inequality 
        \begin{itemize}
            \item Introduction: tbd
            \item Validation Questions: tbd
        \end{itemize}
        \item Tax Filing
        \begin{itemize}
            \item Introduction: tbd
            \item Validation Questions: tbd
        \end{itemize}
        \item Full Info
        \begin{itemize}
            \item Introduction: tbd
            \item Validation Questions: tbd
        \end{itemize}
    \end{itemize}
\end{enumerate}

\section*{General Considerations}
\underline{Considerations on Fairness and Inequality (not related to tax evasion)}
\
\begin{enumerate}[resume]
    \item How serious do you consider the issue of income and wealth inequality to be in your country?\\
    \textit{(1: Not serious at all, 2: Slightly serious, 3: Moderately serious, 4: Serious, 5: Very Serious)}
\end{enumerate}
\begin{enumerate}[resume]
    \item Considering all the responsibilities a government has, how important do you think it is for the government to reduce income and wealth inequality? \\
    \textit{(1: Not important at all, 2: Slightly important, 3: Moderately important, 4: Important, 5: Very important)}
\end{enumerate}
\underline{Considerations on Tax Evasion, Revenue and Involved Inequality}
\begin{enumerate}[resume]
    \item In general, how concerning is the current level of tax evasion in your country?\\
    \textit{(1: Not concerning at all, 2: Slightly concerning, 3: Moderately concerning, 4: Concerning), 5: Very concerning}
\end{enumerate}
\begin{enumerate}[resume]
    \item Do you think tax evasion is mostly driven by a small group of individuals responsible for a large share of total evasion, or do you think many taxpayers contribute roughly equally to the total amount of taxes evaded?\\
    \textit{(1: A small group is responsible for most tax evasion, 2: (no description), 3: (no description), 4: (no description), 5: Many taxpayers contribute equally to total tax evasion)}
    \end{enumerate}
\begin{enumerate}[resume]
    \item Considering all the responsibilities a government has, how important do you think it is for the government to reduce the extent of tax evasion?\\
    \textit{(1: Not important at all, 2: Slightly important, 3: Moderately important, 4: Important, 5: Very important)} 
\end{enumerate}
\underline{Considerations on Tax Filing Costs} \\ 
\begin{enumerate}[resume]
    \item Do you consider tax filing costs in your country to be low or high?\\ 
    \textit{(1: Very low, 2: Low, 3: Neither low nor high, 4: High 5: Very high)}
    \item Considering all the responsibilities a government has, how important do you think it is for the government to reduce the costs associated with filing taxes? \\ \textit{(1: Not important at all, 2: Slightly important, 3: Moderately important, 4: Important, 5: Very important)}
\end{enumerate}
\underline{Considerations on Data Privacy} \\ 
\begin{adjustwidth}{0.75em}{0em}The following question asks about your general view on the federal government. When answering, \uline{please think broadly about the past 10 years and your expectations for the next 10 years, and try to set aside the actions of any specific administration.}\end{adjustwidth}

\begin{enumerate}[resume]
    \item Some people are concerned about governments collecting data on individuals and businesses, viewing it as a potential threat to privacy. Others believe that such data collection is acceptable or even beneficial if it serves public interests, such as ensuring tax compliance or improving public services. \\ How concerned are you about the government collecting data on individuals and businesses?\\
    \textit{(1: Not concerned at all, 2: Slightly concerned, 3: Moderately concerned, 4: Concerned, 5: Very concerned)}
\end{enumerate}
\begin{enumerate}[resume]
    \item Considering all the responsibilities a government has, how important do you think it is for the government to ensure data privacy for its citizens? \\
    \textit{(1: Not important at all, 2: Slightly important, 3: Moderately important, 4: Important, 5: Very important)}
\end{enumerate}    

\section*{Policy Block 1: Information Sharing}
\subsection*{Third Party Reporting: Considerations}
\begin{adjustwidth}{0.75em}{0em}Please carefully read the following explanation of third-party reporting and keep it in mind when answering the following questions.\end{adjustwidth} 
\vspace{0.5em}
\begin{adjustwidth}{1.25em}{1.25em}  \textbf{Explanation:} In order to fulfil its functions, the tax authority uses information from other sources, such as employers, banks and insurance companies. With third-party reporting, this tax-relevant data is automatically sent from domestic firms and institutions to the tax authority. In a comprehensive third-party reporting system, the tax authority receives detailed and automatic reports from many sources in a timely manner.
\\
%\textbf{Explanation:} In order to fulfil its functions, the tax authority uses information from other sources, such as employers, banks and insurance companies. With third-party reporting, this tax-relevant data is automatically sent from domestic firms and institutions to the tax authority. In a comprehensive third-party reporting system, the tax authority receives detailed and automatic reports from many sources.
\end{adjustwidth}
\underline{Considerations on Tax Evasion and Revenue}
\begin{enumerate}[resume]
    \item In your view, how much would comprehensive third-party reporting reduce tax evasion? \\
    \textit{(1: Not at all, 2: Slightly, 3: Moderately, 4: Considerably, 5: Greatly)} 
    \item In your view, how much would comprehensive third-party reporting reduce the administrative costs of tax enforcement for the government?\\
    \textit{(1: Not at all, 2: Slightly, 3: Moderately, 4: Considerably, 5: Greatly)} 
    \item In your view, how much would comprehensive third-party reporting increase government revenue?  \\
    \textit{(1: Not at all, 2: Slightly, 3: Moderately, 4: Considerably, 5: Greatly)} 
\end{enumerate}
\underline{Considerations on Fairness and Inequality}
\begin{enumerate}[resume]
    \item In your view, how much would comprehensive third-party reporting increase the overall fairness of the tax system?\\
    \textit{(1: Not at all, 2: Slightly, 3: Moderately, 4: Considerably, 5: Greatly)} 
    \item In your view, how important is comprehensive third-party reporting for reducing income and wealth inequality?\\
    \textit{(1: Not important at all, 2: Slightly important, 3: Moderately important, 4: Important, 5: Very important)}
    %\item In your view, to what extent does comprehensive third-party reporting reduce the administrative costs of tax enforcement for the government?\\
    %\textit{(1: No reduction at all - 10: Strong reduction)}
\end{enumerate}
\underline{Considerations on Tax Filing Costs}
\begin{enumerate}[resume]    
    \item In your view, how much would comprehensive third-party help in reducing time and monetary costs for taxpayers when filing their taxes?\\
    \textit{(1: Not at all, 2: Slightly, 3: Moderately, 4: Considerably, 5: Greatly)} 
\end{enumerate}
\underline{Considerations on Data Privacy}
\begin{enumerate}[resume]    
    \item In your opinion, how much would comprehensive third-party reporting reduce people's data privacy relative to a tax system with no third-party reporting?\\
    \textit{(1: Not at all, 2: Slightly, 3: Moderately, 4: Considerably, 5: Greatly)}
    \item How concerning would you find a loss of privacy due to comprehensive third-party reporting?\\
    \textit{(1: Not concerning at all, 2: Slightly concerning, 3: Moderately concerning, 4: Concerning, 5: Very concerning)}
\end{enumerate}
\begin{enumerate}[resume]
    \item How concerning a loss of privacy due to third-party reporting is may depend on what data is involved. (In red you can see your response from before regarding your general concern for a loss of privacy due to third-party reporting.) Please indicate how concerning you would find the loss of privacy due to third-party reporting for the following specific types of data:

\begin{enumerate}[label=\alph*)]
        \item General employment data (e.g., job type, hours worked, salary, bonuses)\\
        \textit{(1: Not concerning at all, 2: Slightly concerning, 3: Moderately concerning, 4: Concerning, 5: Very concerning)}
        \item Employment data limited to tax-relevant financial information (e.g., salary, bonuses)\\
        \textit{(1: Not concerning at all, 2: Slightly concerning, 3: Moderately concerning, 4: Concerning, 5: Very concerning)}
        \item General health insurance data (e.g., types of treatments, premium payments, subsidies)\\
        \textit{(1: Not concerning at all, 2: Slightly concerning, 3: Moderately concerning, 4: Concerning, 5: Very concerning)}
        \item Health insurance data limited to tax-relevant financial information (e.g., premium payments, subsidies)\\
        \textit{(1: Not concerning at all, 2: Slightly concerning, 3: Moderately concerning, 4: Concerning, 5: Very concerning)}
        \item General bank account data (e.g., all transactions and balances)\\
        \textit{(1: Not concerning at all, 2: Slightly concerning, 3: Moderately concerning, 4: Concerning, 5: Very concerning)}
        \item Bank account data limited to tax-relevant information (e.g., end-of-year balances, interest payments)\\
       \textit{(1: Not concerning at all, 2: Slightly concerning, 3: Moderately concerning, 4: Concerning, 5: Very concerning)}
        \item General payment data (e.g., credit card or mobile payment histories)\\
        \textit{(1: Not concerning at all, 2: Slightly concerning, 3: Moderately concerning, 4: Concerning, 5: Very concerning)}
        \item Payment data limited to tax-relevant information (e.g., VAT-relevant payment information)\\
        \textit{(1: Not concerning at all, 2: Slightly concerning, 3: Moderately concerning, 4: Concerning, 5: Very concerning)}
\end{enumerate}
\end{enumerate}
\subsection*{Third Party Reporting: Support}
\begin{adjustwidth}{0.75em}{0em}In the following, we will ask you about your support for comprehensive third-party reporting. Below, we repeat the explanation of third-party reporting.\end{adjustwidth}
\vspace{0.5em}
\begin{adjustwidth}{1.25em}{1.25em}  \textbf{Explanation:} In order to fulfill its functions, the tax authority uses information from other sources, such as employers, banks and insurance companies. With third-party reporting, this tax-relevant data is automatically sent from domestic firms and institutions to the tax authority. In a comprehensive third-party reporting system, the tax authority receives detailed and automatic reports from many sources in a timely manner.
\end{adjustwidth}
\begin{enumerate}[resume]
    \item In general, do you support comprehensive third-party reporting of tax-relevant information to the tax authority?  \\ \textit{(1: Strongly Oppose, 2: Oppose, 3: Neither oppose nor support, 4: Support, 5: Strongly support)}
\end{enumerate}

\begin{adjustwidth}{0.75em}{0em}In the following, please state how much your support for comprehensive third-party reporting would change if it contained certain elements or were applied to certain subgroups. (In red you can see the level of support you indicated before.) \end{adjustwidth}  
\begin{enumerate}[resume]
    \item How much do you support comprehensive third-party reporting of tax-relevant information to the tax authority... 

\begin{enumerate}[label=\alph*)]
    \item ...if it applied equally to all taxpayers — both individuals and businesses?\\
    \textit{(1: Strongly Oppose, 2: Oppose, 3: Neither oppose nor support, 4: Support, 5: Strongly support)}
    \item ...if it applied primarily to individual taxpayers, but not to businesses?\\
    \textit{(1: Strongly Oppose, 2: Oppose, 3: Neither oppose nor support, 4: Support, 5: Strongly support)}
    \item ...if it applied primarily to businesses, but not to individual taxpayers?\\
    \textit{(1: Strongly Oppose, 2: Oppose, 3: Neither oppose nor support, 4: Support, 5: Strongly support)}
    \item ...if it focused primarily on high-income individuals, but not on low-income individuals?\\
    \textit{(1: Strongly Oppose, 2: Oppose, 3: Neither oppose nor support, 4: Support, 5: Strongly support)}
    \item ...if it focused primarily on large enterprises, but not on small or medium-sized businesses?\\
    \textit{(1: Strongly Oppose, 2: Oppose, 3: Neither oppose nor support, 4: Support, 5: Strongly support)}
    \item ...if the data could only be used for tax enforcement purposes, and not for other government functions?\\
    \textit{(1: Strongly Oppose, 2: Oppose, 3: Neither oppose nor support, 4: Support, 5: Strongly support)}
    \end{enumerate}
\end{enumerate}    
\subsection*{Third Party Reporting: Petition}
\tk{Q to SS: We have moved the petition directly to the corresponding policy block. The main advantage is that it is much more aligned, and there will probably be less confusion. One potential disadvantage is that some respondents may dislike the petition, which may influence their subsequent answers. What do you think, where should it be placed? \textbf{We suggest to keep it here.}} 
\begin{enumerate}[resume]
    \item Thank you for your answers about comprehensive third party reporting. For interested respondents, we have prepared a related petition that you can support. The petition is in favor of comprehensive third-party reporting to act against tax evasion. As soon as the survey is complete, we will send the results to the [head of state’s] office, informing them what share of people who took this survey were willing to support which petition. \\ \\
    \textbf{Petition for stronger action against tax evasion:} "I agree that fair taxation is essential for a functioning society. The government must act to ensure everyone pays their fair share. One way to reduce tax evasion is through comprehensive third-party reporting---requiring banks, employers, and other institutions to share key financial information with tax authorities in a timely manner. I support expanding such reporting systems to close the tax gap and ensure tax fairness." \\ \\
        Do you want to support this petition? \textit{(Yes / No)} (You will \underline{NOT} be asked to sign, only your answer here is required and remains anonymous.)
\end{enumerate}
\subsection*{International Exchange of Information: Considerations and Support}
\tk{Note to SS: Generally, the structure of questions regarding international exchange of information follows the one about comprehensive third-party reporting (for direct comparability). Otherwise we will note differently.} 

%Please carefully read the following explanation of automatic exchange of information and keep it in mind when answering the following questions.\\ 
%\begin{adjustwidth}{1em}{1em} \textbf{Explanation:} Tax authorities do not only receive information from domestic sources — they also receive data from foreign financial institutions about their residents’ foreign bank accounts. This happens through automatic exchange of information, where countries agree to share bank account data with each other.  
%\end{adjustwidth}
\begin{adjustwidth}{0.75em}{0em} Please carefully read the following explanation of international exchange of information and keep it in mind when answering the following questions.\end{adjustwidth}
\vspace{0.5em}
\begin{adjustwidth}{1em}{1em} \textbf{Explanation:} Tax authorities do not only receive information from domestic sources — they can also obtain data about their residents’ foreign bank accounts through international agreements. In the case of the United States, many foreign banks are required to report information on accounts held by U.S. taxpayers directly or through their local tax authorities to the U.S. government. In some cases, the U.S. also agrees to share similar information with those partner countries, but this depends on the specific agreement in place. \\
\end{adjustwidth}
\underline{Considerations on Tax Evasion and Revenue}
\begin{enumerate}[resume]
    \item In your view, how much does international exchange of information reduce tax evasion? \\
    \textit{(1: Not at all, 2: Slightly, 3: Moderately, 4: Considerably, 5: Greatly)} 
    \item In your view, how much does international exchange of information reduce the administrative costs of tax enforcement for the government?\\
    \textit{(1: Not at all, 2: Slightly, 3: Moderately, 4: Considerably, 5: Greatly)} 
    \item In your view, how much does international exchange of information increase government revenue?  \\
    \textit{(1: Not at all, 2: Slightly, 3: Moderately, 4: Considerably, 5: Greatly)} 
\end{enumerate}
\underline{Considerations on Fairness and Inequality}
\begin{enumerate}[resume]
    \item In your view, how much does international exchange of information increase the overall fairness of the tax system? \\
    \textit{(1: Not at all, 2: Slightly, 3: Moderately, 4: Considerably, 5: Greatly)} 
    \item In your view, how important is international exchange of information for reducing income and wealth inequality?\\
    \textit{(1: Not important at all, 2: Slightly important, 3: Moderately important, 4: Important, 5: Very important)}
\end{enumerate}
\underline{Considerations on Tax Filing Costs}
\begin{enumerate}[resume]
    \item In your view, how much does international exchange of information reduce the time and costs for taxpayers when filing their taxes?\\
    \textit{(1: Not at all, 2: Slightly, 3: Moderately, 4: Considerably, 5: Greatly)}
\end{enumerate}
\underline{Considerations on Data Privacy}
\begin{enumerate}[resume]
    \item In your opinion, how much does international exchange of information with other countries reduce people's data privacy? \\
    \textit{(1: Not at all, 2: Slightly, 3: Moderately, 4: Considerably, 5: Greatly)} 
    \item How concerning would you find a loss of privacy due to international exchange of information? \\ 
    \textit{(1: Not concerning at all, 2: Slightly concerning, 3: Moderately concerning, 4: Concerning, 5: Very concerning)}
\end{enumerate}
\begin{adjustwidth}{0.75em}{0em}\tk{Note to SS: We decided to not add questions regarding data privacy considerations conditional on data types (like we did before for third party reporting).}\end{adjustwidth} 
\underline{Support}
\begin{enumerate}[resume]
    \item Do you support international exchange of information with other countries? \\
    \textit{(1: Strongly Oppose, 2: Oppose, 3: Neither oppose nor support, 4: Support, 5: Strongly support)}
\end{enumerate}
\begin{adjustwidth}{0.75em}{0em} \tk{Note to SS: We decided to not add questions regarding conditional support (like we did before for third party reporting).}\end{adjustwidth} 
\underline{Petition} 
\vspace{0.5em}
\begin{adjustwidth}{0.75em}{0em}\tk{Note to SS: We decided to not add a petition related to international information exchange (like we did before for third party reporting).}\end{adjustwidth}
\section*{Policy Block 2: Tax Software and Prepopulated Tax Returns}
\subsection*{Tax Software: Considerations and Support} 
\vspace{0.5em}
\begin{adjustwidth}{0.75em}{0em}Please carefully read the following explanation of tax software and keep it in mind when answering the following questions. \end{adjustwidth}
\vspace{0.5em}
\begin{adjustwidth}{1.25em}{1.25em} \textbf{Explanation:} Tax software is a computer program used for filing tax returns. It usually provides guidance through tax formulas and explains which deductions might apply..
\end{adjustwidth}
\vspace{0.5em}
\underline{Considerations on Tax Filing Costs}
\begin{enumerate}[resume]
    \item Assume you obtain the tax software for free. In your view, how much does tax software reduce time and monetary costs for taxpayers when filing their taxes? \\
    \textit{(1: Not at all, 2: Slightly, 3: Moderately, 4: Considerably, 5: Greatly)} 
    \item In some countries, the government provides free tax software as a public service. In others, taxpayers can purchase software from private providers. In terms of reducing time and monetary costs for taxpayers, do you think it is better if private companies or the tax authority provides tax software? \\
    \textit{(1: Definitely private companies, 2: Probably private companies, 3: Both equally, 4: Probably the tax authority, 5: Definitely the tax authority)}
\end{enumerate}
\underline{Considerations on Tax Revenue (Government expenses)}
\begin{enumerate}[resume]
    \item How costly do you think providing tax software would be for the U.S. government? \\
    \textit{(1: Not costly at all, 2: Slightly costly, 3: Moderately costly, 4: Costly, 5: Very costly)} 
    \item In terms of overall costs to the society, do you think it is more efficient if private companies or the tax authority provide tax software?  \\
    \textit{(1: Definitely private companies, 2: Probably private companies, 3: Both equally, 4: Probably the tax authority, 5: Definitely the tax authority)}
\end{enumerate}
\underline{Considerations on Data Privacy}
\begin{enumerate}[resume]
    \item When filing taxes with tax software, taxpayers often need to share private data. How problematic do you think this might be in terms of data privacy? \\
    \textit{(1: Not problematic at all, 2: Slightly problematic, 3: Moderately problematic, 4: Problematic, 5: Very problematic)}
    \item In terms of data privacy, do you think it is better if private companies or the tax authority provide tax software and, therefore, manage tax relevant data? \\
    \textit{(1: Definitely private companies, 2: Probably private companies, 3: Both equally, 4: Probably the tax authority, 5: Definitely the tax authority)}
\end{enumerate}
\underline{Considerations on Fairness and Inequality}
\vspace{0.5em}
\begin{adjustwidth}{0.75em}{0em} \tk{Note to SS: We suggest to not add questions regarding fairness and inequality.}\end{adjustwidth} 
\vspace{0.5em}
\underline{General Trade Off Private vs. Government Provider}
\begin{enumerate}[resume]
    \item Generally, do you think it is better if private companies or the tax authority provides tax software? \\
    \textit{(1: Definitely private companies, 2: Probably private companies, 3: Both equally, 4: Probably the tax authority, 5: Definitely the tax authority)}
\end{enumerate}
\underline{Support} \\
\begin{enumerate}[resume]
    \item Would you support it if the government provided free tax software?\\
    \textit{(1: Strongly Oppose, 2: Oppose, 3: Neither oppose nor support, 4: Support, 5: Strongly support)}
\end{enumerate}
\subsection*{Prepopulated Tax Returns: Considerations and Support} 
\tk{Note to SS: This block generally follows the same structure as the one for tax software, but now we do not ask for the government vs. private provider trade-off. }
\vspace{0.5em}
\begin{adjustwidth}{0.75em}{0em}Please carefully read the following explanation of prepopulated tax returns and keep it in mind when answering the following questions. \end{adjustwidth}
\vspace{0.5em}
\begin{adjustwidth}{1.25em}{1.25em} \textbf{Explanation:} Prepopulated tax returns are forms automatically filled in by the tax authority  based on data they already have (for example data on salaries or interest payments). Taxpayers can review the information, make corrections or additions, and then submit the return.
\end{adjustwidth}
\vspace{0.5em}
\underline{Considerations on Tax Filing Costs}
\begin{enumerate}[resume]
    \item How much do you think prepopulated tax returns reduce time and monetary costs for taxpayers when filing their taxes?\\
    \textit{(1: Not at all, 2: Slightly, 3: Moderately, 4: Considerably, 5: Greatly)}
\end{enumerate}
\underline{Considerations on Tax Revenue (Government expenses)}
\begin{enumerate}[resume]
     \item How costly do you think prepopulating tax returns would be for the U.S. government? \\
    \textit{(1: Not costly at all, 2: Slightly costly, 3: Moderately costly, 4: Costly, 5: Very costly)} 
\end{enumerate}
\underline{Considerations on Data Privacy}
\begin{enumerate}[resume]
     \item To prepopulate tax returns the tax authority needs a significant amount of tax relevant data. How problematic do you think this might be in terms of data privacy? \\ 
     \textit{(1: Not problematic at all, 2: Slightly problematic, 3: Moderately problematic, 4: Problematic, 5: Very problematic)} 
\end{enumerate}
\underline{Considerations on Fairness and Inequality}
\vspace{0.5em}
\begin{adjustwidth}{0.75em}{0em} \tk{Note to SS: We decided to not add questions regarding fairness and inequality.} \end{adjustwidth} 
\underline{Support}
\begin{enumerate}[resume]
    \item Would you support it if the government offered free prepopulated tax returns?\\
    \textit{(1: Strongly Oppose, 2: Oppose, 3: Neither oppose nor support, 4: Support, 5: Strongly support)}
\end{enumerate}
\underline{Updated Support for Third Party Reporting} \\
\begin{enumerate}[resume]
    \item Comprehensive third-party reporting is a key requirement for governments to offer prepopulated tax returns, because otherwise the government cannot prefill the returns. Earlier, you were asked about your support for comprehensive third-party reporting in general (in red you can see the level of support you indicated before). Now, please state how much you would support comprehensive third-party reporting of tax-relevant information to the tax authority... 
    \begin{enumerate}
    \item ...if you knew that the collected data would be used to offer prepopulated tax returns (in addition to tax enforcement purposes).\\ \textit{(1: Strongly Oppose, 2: Oppose, 3: Neither oppose nor support, 4: Support, 5: Strongly support)}  
    \item ...if the data were \underline{only} used to provide prepopulated tax returns (and not for enforcement purposes).\\ \textit{(1: Strongly Oppose, 2: Oppose, 3: Neither oppose nor support, 4: Support, 5: Strongly support)}
    \item ...if taxpayers could voluntarily choose to have their own data covered by third-party reporting, which would allow the tax authority to prepopulate tax returns. \\
    \textit{(1: Strongly Oppose, 2: Oppose, 3: Neither oppose nor support, 4: Support, 5: Strongly support)}
\end{enumerate}
\end{enumerate}
\underline{Voluntary Disclosure of Data} \\
\begin{enumerate}[resume]
    \item Closely related, if taxpayers could voluntarily choose to have their own data covered by third-party reporting, would you be willing to have your tax-relevant data - for example your bank account data - covered by automatic reporting to the tax authority, ...?   \\
    \begin{enumerate}
    \item ...if you knew that the collected data would \underline{only} be used for tax enforcement  purposes (and not for prepopulating tax returns).\\     \textit{(1: Not willing at all, 2: Slightly willing, 3: Moderately willing, 4: Willing, 5: Very willing)}
    \item ...if you knew that the collected data would be used to offer prepopulated tax returns (in addition to tax enforcement purposes).\\     \textit{(1: Not willing at all, 2: Slightly willing, 3: Moderately willing, 4: Willing, 5: Very willing)}
    \item ...if the data were \underline{only} used to provide prepopulated tax returns (and not for enforcement purposes).\\     \textit{(1: Not willing at all, 2: Slightly willing, 3: Moderately willing, 4: Willing, 5: Very willing)}
\end{enumerate}
\end{enumerate}
\subsection*{Tax Software and Prepopulated Tax Returns: Petition} 
\begin{enumerate}[resume]
    \item Thank you for your answers about tax software and prepopulated tax returns. For interested respondents, we have prepared a related petition that you can support. The petition is in favor of making tax filing easier and more accessible for everyone. As soon as the survey is complete, we will send the results to the [head of state’s] office, informing them what share of people who took this survey were willing to support which petition. \\ \\
    \textbf{Petition to make tax filing easier and more accessible for everyone:} "I believe that paying taxes should be simple and transparent. Many governments already have the information needed to prepare tax returns for most taxpayers. I support the introduction of prepopulated tax returns and government-provided tax filing services, such as free tax software, to save time, reduce errors, and make tax compliance easier for everyone."\\ \\
    Do you want to support this petition? \textit{(Yes / No)} (You will \underline{NOT} be asked to sign, only your answer here is required and remains anonymous.)
\end{enumerate}
\section*{Final Questions}
\tk{Q to SS: \textbf{We suggest to not ask the following question.} The advantage of being able to obtain an "unbiased" result might be smaller than the disadvantage of the risk of obtaining a "biased" result. However, as you used it in previous surveys, we wanted to ask whether you think it should be included?}
\begin{enumerate}[resume]
    \stepcounter{enumi}
    \item[\sout{\theenumi.}] \textcolor{gray}{Do you think this survey was biased?\\ \textit{(Yes, left-wing biased / Yes, right-wing biased / No)}}
    \item Please feel free to give us any feedback or impressions regarding this survey.
\end{enumerate}    
%\section*{Petition}
%\subsection*{Support for Comprehensive Third-Party Reporting}
%\begin{enumerate}[resume]
%    \item Thanks for participating in the survey. For interested respondents we have prepared three related petitions that you can support: 
%    \begin{itemize}
%        \item Petition a) is in favor comprehensive third-party reporting to act against tax evasion. 
%        \item Petition b) is against comprehensive third-party reporting in order to protect privacy and government overreach in the tax system. 
%        \item Petition c) is in favor of making tax filing easier and more accessible for everyone. 
%    \end{itemize}
%    As soon as the survey is complete, we will send the results to the [head of state’s] office, informing them what share of people who took this survey were willing to support which petition. \\ \\ Would you like to have a look at the petitions to see if you want to support one or more of them? (You will \underline{NOT} be asked to sign, only your answer here is required and remains anonymous.)\\  \textit{(Yes, I want to look at them / No, I do not want to look at them)} 
%     
%    \begin{enumerate}[label=\alph*)]
%        \item \textbf{Petition for stronger action against tax evasion:} "I agree that fair taxation is essential for a functioning society. The government must act to ensure everyone pays their fair share. One way to reduce tax evasion is through comprehensive third-party reporting---requiring banks, employers, and other institutions to share key financial information with tax authorities. I support expanding such reporting systems to close the tax gap and ensure tax fairness." \\ \\
%        Do you want to support this petition? \textit{(Yes / No)} (You will \underline{NOT} be asked to sign, only your answer here is required and remains anonymous.)\\
%        \item \textbf{Petition to protect privacy and limit government overreach in the tax system:} "I believe that while reducing tax evasion is important, it must not come at the expense of individual privacy and personal freedom. Comprehensive third-party reporting creates data collection systems that can be vulnerable to misuse or breaches. Expanding such reporting risks turns financial institutions into surveillance tools for the government. I oppose further expansion of third-party reporting and support a tax system that respects privacy and individual rights."\\ \\
%        Do you want to support this petition? \textit{(Yes / No)} (You will \underline{NOT} be asked to sign, only your answer here is required and remains anonymous.) \\
%        \item \textbf{Petition to make tax filing easier and more accessible for everyone:} "I believe that paying taxes should be simple and transparent. Many governments already have the information needed to prepare tax returns for most taxpayers. I support the introduction of prepopulated tax returns or government-provided tax filing services, such as free tax software, to save time, reduce errors, and make tax compliance easier for everyone."\\ \\
%        Do you want to support this petition? \textit{(Yes / No)} (You will \underline{NOT} be asked to sign, only your answer here is required and remains anonymous.)
%    \end{enumerate}
%\end{enumerate}   
\end{document}